\documentclass{article}
\usepackage{etoolbox}
\usepackage{graphicx}
\usepackage{amsmath}

\ifdef{\myoutput}{}{\def\myoutput{concept}}
\ifdef{\myversion}{}{\def\myversion{1}}

\usepackage[output=\myoutput
,numberofversions=4
,version=\myversion
,seed=4
,randomizequestions=true
,randomizeanswers=true
,writeRfile=true
]{mcexam}

\usepackage[margin=0.95in]{geometry}
\usepackage{multicol}

\usepackage{fancyhdr,lastpage}
\pagestyle{fancy}
\fancyhf{}
\renewcommand{\headrulewidth}{1pt} 
\renewcommand{\footrulewidth}{1pt}
\lfoot{\mctheversion}
\rfoot{Page \thepage\ of \pageref{LastPage}}
\lhead{Quiz 9.2-3 and 10.2 \mctheversion \\ MAT 2000: Introduction to Statistics}
\rhead{Name Print): \hspace{2.5 in}}

\usepackage{calc}
\renewenvironment{setmcquestion}
{\begin{minipage}[t]{\linewidth-\labelwidth}}
	{\end{minipage}\par}
	
\renewenvironment{setmcquestioninstruction}
{\begin{minipage}{\textwidth}}
	{\end{minipage}}
	
\renewenvironment{setmcanswers}{}{}
\setlist[setmcquestions]{label=\mcquestionlabelfmt{*}.
	,ref=\mcquestionlabelfmt{*}
	,itemsep=0.5\baselineskip
	,topsep=1\baselineskip
}

\renewcommand\mcversionlabelfmt[1]{\arabic{#1}}

%copy-pasting from MS Word, use regex to find [\*]*[a-d]\. and replace with \answer 

\raggedbottom

\begin{document}
	Circle the correct answer for each question.  If you're unsure, you can write numbers 1, 2, 3, etc. next to the choices to get partial credit.  Questions refer to the nonindented paragraph above unless otherwise stated.\\[.5cm]
	
	\begin{mcquestions}
		
		\question Suppose the chi-square statistic for a test was 6.991. If instead, the chi-square statistic had been 5.332, would the p-value be smaller, larger, or remain the same?
		\begin{mcanswerslist}
			\answer Smaller
			\answer[correct] Larger
			\answer Remain the same
			\answer Cannot be determined
		\end{mcanswerslist}
		
		\begin{mcquestioninstruction}
			Some student researchers want to determine if students yawning (when being yawned at) is associated with the time of day (morning, afternoon, and evening). The students found the following results from yawning in front of people around campus:
			\begin{itemize}
				\item Morning: 25 of the 39 people yawned back
				\item Afternoon: 15 of the 38 people yawned back
				\item Evening: 29 of the 44 people yawned back
			\end{itemize}
		\end{mcquestioninstruction}
		\question When this test is run, a p-value of 0.03 is found. What would be an appropriate conclusion?
		\begin{mcanswerslist}
			\answer Strong evidence exists that the probability a person yawns is different between all three times of day: morning, afternoon, and evening.
			\answer[correct] Strong evidence exists of an association between time of day and whether a person yawns.
			\answer Strong evidence of no association between time of day and whether a person yawns.
			\answer Strong evidence does not exist of an association between time of day and whether a person yawns.
		\end{mcanswerslist}
		
		\question[follow] After a chi-square test, follow-up confidence intervals for all pairwise differences in proportions are found as follows:
		\begin{itemize}
			\item 0.0299 to 0.4627:  Morning – Afternoon
			\item -0.2237 to 0.1876:  Morning – Evening
			\item -0.4736 to -0.0551:  Afternoon – Evening
		\end{itemize}
		Based off of these intervals, which of the following conclusions are correct?
		\begin{mcanswerslist}
			\answer The probability of yawning in the morning differs significantly from the probability of yawning in the evening.
			\answer[correct] The probability of yawning in the morning differs significantly from the probability of yawning the afternoon.
			\answer The probability of yawning in the afternoon does not differ significantly from the probability of yawning in the evening.
			\answer The probability of yawning in the morning does not differ significantly from the probability of yawning in the afternoon.
		\end{mcanswerslist}
		
		
		\question 3) Why do we do overall tests when comparing multiple proportions and not just do the follow-up confidence intervals?
		\begin{mcanswerslist}
			\answer[correct] Doing an overall test will more likely lead to significant results than just doing the follow-up confidence intervals.
			\answer Doing an overall test allows us to see exactly which group is significantly different from which other groups.
			\answer Doing an overall test allows us to quickly get our results, while the follow-up confidence intervals are very time-consuming.
			\answer Doing an overall test allows us to keep the probability of a type I error at 5\% or whatever significance level we would like.
		\end{mcanswerslist}
		
		
		\question Can telling a joke affect whether or not a customer in a coffee bar leaves a tip for the waiter? A study investigated this question at a coffee bar in France. The waiter randomly assigned coffee-ordering customers into three groups: one received a card telling a joke with their bill, another group received a card containing an advertisement for a local restaurant, and the third received no card. The results are shown below.
		\begin{center}
			\begin{tabular}{| l | l | l | l | l |}
				\hline
				& \textbf{Joke Card} & \textbf{Ad Card} & \textbf{No Card} & \textbf{Total} \\
				\hline
				\textbf{Left Tip} & 32 (0.40) & 16 (0.20) & 24 (0.30) & 72 \\
				\hline
				\textbf{No Tip} & 48 & 64 & 56 & 168 \\
				\hline
				\textbf{Total} & 80 & 80 & 80 & 240 \\
				\hline
			\end{tabular}
		\end{center}
		Are validity conditions met to conduct a chi-square test of association?
		\begin{mcanswerslist}
			\answer Yes, since 80 is larger than 20.
			\answer Yes, since 240 is larger than 20.
			\answer[correct] Yes, since each cell in the table has at least 10 observations.
			\answer No, the sample is not a random sample of customers.
		\end{mcanswerslist}
		
		\begin{mcquestioninstruction}
			According to the 2019 U.S. Census Bureau estimates, 63.4\% of the U.S. self-identifies as Non-Hispanic white, 15.3\% as Hispanic and Latino, 13.4\% as Black or African American, 5.9\% as Asian, and 2.0\% as Other. A random sample of 1000 students graduating from California colleges and universities resulted in the following data on self-identified race:
			\begin{center}
				\begin{tabular}{| p{1in} | p{1in} | p{1in} | p{1in} | p{1in} | p{1in} |}
					\hline
					\textbf{Self-identified Race} & Non-Hispanic White & Black or African American & Hispanic or Latino & Asian & Other \\
					\hline
					\textbf{Number in Sample} & 679 & 51 & 77 & 190 & 3 \\
					\hline
				\end{tabular}
			\end{center}
			Use this information to answer questions \ref{census-1} and \ref{census-2}
		\end{mcquestioninstruction}
		
		\question In the null hypothesis, what does the symbol \(\mu_{\text{Race}}\) represent? \label{census-1}
		\begin{mcanswerslist}
			\answer  The proportion of the U.S. that self-identify as a certain Race
			\answer  The proportion of Californians that self-identify as a certain Race
			\answer[correct]  The proportion of all students graduating from California colleges and universities that self-identify as a certain Race
			\answer  The proportion of the sample of 1000 students that self-identified as a certain Race
		\end{mcanswerslist}
		
		\question[follow] Based upon a p-value of less than 0.001, what is an appropriate conclusion for the goodness-of-fit test? \label{census-2}
		\begin{mcanswerslist}
			\answer  Strong evidence does not exist that the distribution of self-identified race among California college and university students differs from the distribution of self-identified race in the U.S. overall.
			\answer[correct]  Strong evidence exists that the distribution of self-identified race among California college and university students differs from the distribution of self-identified race in the U.S. overall.
			\answer  Strong evidence exists that the distribution of self-identified race among California college and university students is the same as the distribution of self-identified race in the U.S. overall.
			\answer  Strong evidence exists that the probabilities of each self-identified race among California college and university students are not all equal.
		\end{mcanswerslist}
		
		\begin{mcquestioninstruction}
			Are people equally likely to be born on any day of the seven days of the week? Or are some days more likely to be a person's birthday than other days?  To investigate this question, days of birth were recorded for the 147 "noted writers of the present" listed in The World Almanac and Book of Facts 2000. The counts for the seven days of the week are given in the following table.
			\begin{center}
				\begin{tabular}{ccccccc}
					Mon&Tues&Wed&Thurs&Fri&Sat&Sun\\ \hline
					17&26&22&23&18&15&25					
				\end{tabular}
			\end{center}
			Use this information to answer questions \ref{birthday-1} and \ref{birthday-2}
		\end{mcquestioninstruction}
		\question Are validity conditions met to conduct a theory-based chi-square goodness-of-fit test? \label{birthday-1}
		\begin{mcanswerslist}
			\answer[correct]  Yes, since all observed counts are at least 10.
			\answer  Yes, since we have at least 10 successes and 10 failures.
			\answer  No, since the p-values in questions 35 and 37 were different.
			\answer  No, since this was not a random sample.
		\end{mcanswerslist}
		
		\question[follow] Based upon a p-value of 0.53, what is the appropriate conclusion for the test? \label{birthday-2}
		\begin{mcanswerslist}
			\answer  Strong evidence exists that people are equally likely to be born on any day of the seven days of the week.
			\answer  Strong evidence does not exist that people are equally likely to be born on any day of the seven days of the week.
			\answer  Strong evidence exists that people are not equally likely to be born on any day of the seven days of the week.
			\answer[correct]  Strong evidence does not exist that people are not equally likely to be born on any day of the seven days of the week.
		\end{mcanswerslist}
		
		\question The two graphs, A and B, show dotplots from two different data sets. Suppose I want to compare the means in A and do another comparison of means in B. For each case an F-statistic is computed. Which of the following will be true about these statistics?
		\begin{center}
			\includegraphics[width = 0.75\textwidth]{../images/q-10-2-1.png}
		\end{center}
		\begin{mcanswerslist}
			\answer The F-statistic is the same in A and B.
			\answer The F-statistic is larger in A than in B.
			\answer[correct] The F-statistic is smaller in A than in B.
			\answer We cannot tell if the F-statistic will be larger in A or B.
		\end{mcanswerslist}
		
		\question An experiment is conducted which randomly places students into one of three treatment groups: (1) those students who will take a test that is printed on white paper, (2) those students who will take the same test, but printed on red paper, and (3) those who will take the same test, but printed on blue paper. The researcher analyzes the data using ANOVA in order to compare the average scores in the three groups looking for evidence that the average score is different on at least one of the three tests. Assume that at least one group really is different when answering the following question.
		
		As the size of the sample in each group increases, all else remaining the same, what will be the impact on the strength of evidence for finding a difference in the average scores on the three exams?
		
		\begin{mcanswerslist}
			\answer[correct] There will be more evidence of a difference.
			\answer There will be less evidence of a difference.
			\answer There will be no impact on the strength of evidence of a difference.
			\answer There is not enough information to tell.
		\end{mcanswerslist}
		
		\begin{mcquestioninstruction}
			A British study examined whether the type of background music playing in a restaurant affected the amount of money that diners spent on their meals.  The researchers asked a restaurant to alternate silence, popular music, and classical music on successive nights over eighteen days.  Each type of music was played for 6 nights (the order was randomly determined to guard against confounding). The type of music played and the amount spent on food and drinks (in British pounds) for each of 393 customers was recorded. The data are summarized below.
			\begin{center}
				\begin{multicols}{2}
					\begin{tabular}{l|c|c|c|}
						&Sample size&Mean&StDev\\
						Classical&120&24.1&2.30\\
						Pop&142&21.9&2.73\\
						None&131&21.7&3.38\\
						Pooled&393&22.5&2.85
					\end{tabular}
					
					\includegraphics[width=0.5\textwidth]{../images/q-10-2-2.png}
				\end{multicols}
			\end{center}
		\end{mcquestioninstruction}
		\question All of the following statements are true for this situation, but one is not a legitimate validity condition.  Which one is not the real validity condition?
		\begin{mcanswerslist}
			\answer The standard deviations of the sample are approximately equal.
			\answer The sample sizes are all at least 20.
			\answer The sample distributions show no strong skewness.
			\answer[correct] The sample means are all at least 20.
		\end{mcanswerslist}
		
		\question[follow] The p-value for the ANOVA F-test is less than 0.001. How would you interpret this value?
		\begin{mcanswerslist}
			\answer[correct] In less than 0.1\% of all samples, the F-statistic would be as large or larger than the one observed, assuming there is no association between music type and amount spent on food and drinks.
			\answer In less than 0.1\% of all samples, the F-statistic would be as small or smaller than the one observed, assuming there is no association between music type and amount spent on food and drinks.
			\answer In less than 0.1\% of all samples, the F-statistic would be as large or larger than the one observed, assuming there is an association between music type and amount spent on food and drinks.
			\answer In less than 0.1\% of all samples, the F-statistic would be as small or smaller than the one observed, assuming there is an association between music type and amount spent on food and drinks.
		\end{mcanswerslist}
		
		\question When performing an ANOVA F-test, which definition of a p-value is the most accurate?
		\begin{mcanswerslist}
			\answer The p-value is the probability that the observed F-statistic will occur again.
			\answer[correct] If the null hypothesis is true, the p-value is the probability of observing an F-statistic as large or larger than the one actually observed.
			\answer The p-value is the value that an F-statistic must reach in order to be significant under the null hypothesis.
			\answer The p-value is the probability that the null hypothesis is true.
		\end{mcanswerslist}
	\end{mcquestions}
\end{document}